% Problem record op_5083d02a18761d8a. This TeX file is derived from
% database/problems_json/op_5083d02a18761d8a.json by scripts/sync-tex.mjs; edit the JSON record
% and rerun the script rather than editing this file.
\section{Universality of arbitrary self-adjoint polynomial Hamiltonians}
\label{sec:op_5083d02a18761d8a}

\paragraph{Problem.}
Can every physical bosonic unitary be approximated by a single self-adjoint polynomial Hamiltonian evolution at any finite input energy?
Fix $n\geq1$ and a unitary $U$ on $L^2(\mathbb R^n)$ with $U\mathcal S(\mathbb R^n)\subseteq\mathcal S(\mathbb R^n)$.
Here $\mathcal S$ is the Schwartz space of smooth, rapidly decreasing functions.
The canonical operators obey $[q_j,p_k]=i\delta_{jk}$.
Define the total photon number and the energy-constrained channel distance by
\begin{equation}
 \begin{aligned}
 N&=\frac12\sum_{j=1}^n(q_j^2+p_j^2-1),\\
 \|\mathcal U-\mathcal V\|_{\diamond,E}
 &=\sup_{\rho_{AR}:\operatorname{Tr}(\rho_A N)\leq E}
 \|[(\mathcal U-\mathcal V)\otimes\operatorname{id}_R](\rho_{AR})\|_1 .
 \end{aligned}
 \label{eq:poly-universal-distance}
\end{equation}
In Eq.~\eqref{eq:poly-universal-distance}, $\mathcal U(\rho)=U\rho U^\dagger$ and $R$ is any auxiliary reference system.
For every finite $E\geq0$ and $\varepsilon>0$, does there exist a real Weyl-ordered polynomial $P(q_1,p_1,\ldots,q_n,p_n)$ with a self-adjoint realization such that
\begin{equation}
 \|\mathcal U-\mathcal V_P\|_{\diamond,E}<\varepsilon,
 \qquad
 \mathcal V_P(\rho)=e^{-iP}\rho e^{iP}?
 \label{eq:poly-universal-target}
\end{equation}
The polynomial and its degree in Eq.~\eqref{eq:poly-universal-target} may depend on $U,E,\varepsilon$.

\subsection*{Status}

Solved

\subsection*{Source}

Arzani, Booth, and Chabaud establish polynomial universality in Theorems~1--3, with the finite-mode finite-block construction in arXiv Appendix~B \sourcecite{ref:poly-universal-abc}{ABC25}.
The self-adjoint realization required here is supplied by the explicit completion below.
That completion is an editorial deduction, not a claim that the paper proves self-adjointness of its unmodified polynomial.

\subsection*{Progress}

\begin{itemize}
  \item Theorem~2 constructs a symmetric polynomial realizing any Hermitian matrix on a finite Fock block.
The block is invariant.
Theorems~1 and 3 combine finite-dimensional unitary approximation with this construction \sourcecite{ref:poly-universal-abc}{ABC25}.

  \item To complete the domain argument, let $B$ be a symmetric polynomial of degree $d$ realizing a chosen finite block.
Let $K$ bound the total photon numbers in that block.
On the finite Fock span,
$\|B\psi\|\leq C\|(N+1)^{d/2}\psi\|$.
Choose an integer $r$ with $2r(K+1)>d/2$ and set
\begin{equation}
 Q(N)=\left[\prod_{j=0}^{K}(N-j)\right]^{2r},
 \qquad P=B+Q(N).
 \label{eq:poly-universal-completion}
\end{equation}
The operator $Q(N)$ in Eq.~\eqref{eq:poly-universal-completion} is self-adjoint on its spectral domain.
It vanishes on the chosen block and dominates $(N+1)^{d/2}$ at large photon number.
Thus $B$ has relative bound zero with respect to $Q(N)$.
The Kato--Rellich theorem gives a self-adjoint closure of $P$ with the same finite block \sourcecite{ref:poly-universal-teschl}{Tes09}.

  \item For a cutoff at total photon number $L$, the discarded input probability is at most $E/(L+1)$.
This bound also holds with a reference system.
Finite-dimensional approximation of the images of the retained basis vectors gives an approximating unitary on a larger finite Fock block.
The completion in Eq.~\eqref{eq:poly-universal-completion} realizes that block exactly.
Taking $L$ and then the output block large enough proves Eq.~\eqref{eq:poly-universal-target} for every finite number of modes.
\end{itemize}

\subsection*{References}

\begin{enumerate}[label={},leftmargin=4.8em,labelsep=0.5em]
  \footnotesize
  \item[\textup{[ABC25]}]\label{ref:poly-universal-abc}
  F. Arzani, R. I. Booth, and U. Chabaud, "Effective Descriptions of Bosonic Systems Can Be Considered Complete," \emph{Nature Communications} \textbf{16}, 9744 (2025). \href{https://doi.org/10.1038/s41467-025-64872-3}{doi:10.1038/s41467-025-64872-3}; \href{https://arxiv.org/abs/2501.13857}{arXiv:2501.13857}.

  \item[\textup{[Tes09]}]\label{ref:poly-universal-teschl}
  G. Teschl, \emph{Mathematical Methods in Quantum Mechanics: With Applications to Schr\"odinger Operators}, Graduate Studies in Mathematics \textbf{99}, American Mathematical Society (2009), Theorem~6.4, p.~135. \href{https://www.mat.univie.ac.at/~gerald/ftp/book-schroe/schroe.pdf}{Author-hosted full text}.
\end{enumerate}

\subsection*{Comment}

The finite-block universality and approximation results are peer-reviewed.
The self-adjoint completion above applies the standard Kato--Rellich theorem to remove the complement-domain issue.
It can increase the polynomial degree, so no degree-$3d$ claim is made for the completed Hamiltonian.
This result permits a different polynomial for each target and accuracy.
It does not prove universality of an arbitrary fixed non-Gaussian generator with Gaussian controls.

\subsection*{Field}

Quantum algorithm

\subsection*{Topic}

Quantum circuit complexity; Hamiltonian simulation

\subsection*{ID}

\texttt{op\_5083d02a18761d8a}
